% ============================================================
\section{Priorización de los drivers}\label{sec:priorizacion}
% ============================================================

El orden de las iteraciones no es el orden en que se numeraron los drivers, sino el del impacto que cada uno tiene sobre la arquitectura. Se atiende primero lo que decide la forma del sistema y lo que sería más caro de cambiar después, porque una iteración posterior puede acomodarse a una decisión ya tomada, pero no al revés.

% ------------------------------------------------------------
\subsection{Criterios}

Cada driver se califica con tres criterios, de 1 a 5. El total ordena las iteraciones.

\begin{tabla}{|L{3.4cm}|Y|}
\hline
\filacab \cab{Criterio} & \cab{Qué mide} \\ \hline
\endhead
Alcance estructural &
Cuántos elementos del sistema cambian si el driver se resuelve de otra manera. Un 5 significa que cambian a la vez cliente, servidor, protocolo y capa de datos; un 1, que el efecto queda dentro de un componente. \\ \hline
Costo de cambio tardío &
Qué cuesta cambiar la decisión cuando ya hay código construido. Pesa aquí que haya clientes desplegados, datos escritos con un formato o componentes construidos sobre el supuesto anterior. \\ \hline
Drivers que habilita &
Cuántos de los demás drivers dependen de que este ya esté resuelto para poder decidirse. Un driver que no habilita a ninguno puede esperar sin bloquear el avance. \\ \hline
\end{tabla}

% ------------------------------------------------------------
\subsection{Calificación y orden de las iteraciones}

\begin{tabla}{|L{1.4cm}|L{4.5cm}|L{1.4cm}|L{1.4cm}|L{1.4cm}|L{1.0cm}|L{1.3cm}|}
\hline
\filacab \cab{Driver} & \cab{Enunciado abreviado} & \cab{Alcance} & \cab{Costo tardío} & \cab{Habilita} & \cab{Total} & \cab{Iteración} \\ \hline
\endhead
DRV-01 & Servidor autoritativo sobre el estado de la partida & 5 & 5 & 5 & 15 & 1 \\ \hline
DRV-02 & Protocolo propio sobre TCP & 4 & 5 & 4 & 13 & 2 \\ \hline
DRV-05 & Un único motor de reglas, parametrizable y sin interfaz & 4 & 4 & 4 & 12 & 3 \\ \hline
DRV-03 & Presupuesto de un segundo por jugada & 5 & 4 & 3 & 12 & 4 \\ \hline
DRV-04 & Política explícita de persistencia por tipo de dato & 3 & 4 & 3 & 10 & 5 \\ \hline
DRV-06 & Protección del menor: control en el servidor y datos mínimos & 3 & 3 & 2 & 8 & 6 \\ \hline
DRV-07 & Localización de extremo a extremo & 3 & 4 & 1 & 8 & 7 \\ \hline
DRV-08 & Plazo de 14 semanas: rebanadas verticales y puntos de variación & 2 & 2 & 2 & 6 & 8 \\ \hline
\end{tabla}

\textbf{Cómo se rompen los empates.} DRV-05 y DRV-03 empatan en 12 y se ordenan por dependencia: DRV-03 solo puede repartir el segundo entre validación, red y persistencia si antes existe una sola implementación del motor, porque validar en el cliente sin un motor compartido produce reversiones y rompe la medida de ESC-01. DRV-06 y DRV-07 empatan en 8 y se ordenan por la fuerza de su origen: DRV-06 nace de cuatro requisitos impuestos y de stakeholders con poder de veto (STK-02, STK-17), mientras que DRV-07 nace de un objetivo de negocio y de convenciones que ningún componente incumple por su cuenta si se fijan a tiempo.

\textbf{DRV-08 se aplica además de forma transversal.} Es el último en tener iteración propia porque no cambia qué componentes existen, pero su exigencia de rebanadas verticales gobierna el orden en que se construye lo que las demás iteraciones deciden, y sus tres puntos de variación se verifican en cada incremento semanal. Su iteración se ocupa de acotarlos, no de descubrirlos.

% ------------------------------------------------------------
\subsection{Plan de iteraciones}\label{sec:priorizacion:plan}

Este documento desarrolla la primera iteración de cada driver, en el orden de la tabla anterior. Cada una tiene un solo objetivo, un solo elemento y una sola medida, y termina en un incremento demostrable.

\begin{tabla}{|L{1.2cm}|L{1.5cm}|L{4.6cm}|Y|}
\hline
\filacab \cab{Iter.} & \cab{Driver} & \cab{Objetivo único} & \cab{Elemento a descomponer y medida} \\ \hline
\endhead
1 & DRV-01 & El servidor es la única autoridad sobre el estado de una partida en curso. & Gestión del estado de la partida. Medida: ESC-05. \\ \hline
2 & DRV-02 & El tráfico capturado no revela nada legible ni se puede reutilizar. & Servicio de protocolo. Medida: ESC-06. \\ \hline
3 & DRV-05 & Las reglas cambian sin tocar código y el motor se ejecuta sin interfaz. & Motor de reglas. Medida: ESC-12. \\ \hline
4 & DRV-03 & Una jugada se responde en menos de un segundo. & Camino de la jugada. Medida: ESC-01. \\ \hline
5 & DRV-04 & Una caída de la base de datos pierde solo lo acordado y deja constancia. & Servicio de acceso a datos. Medida: ESC-09. \\ \hline
6 & DRV-06 & Ningún mensaje inapropiado llega a un menor, ni siquiera cuando el filtro falla. & Servicio de chat. Medida: ESC-04. \\ \hline
7 & DRV-07 & Se agrega un idioma nuevo sin tocar código ni rehacer pantallas. & Catálogo de textos y formatos. Medida: ESC-11. \\ \hline
8 & DRV-08 & El oponente controlado por IA se puede recortar sin tocar el núcleo. & Punto de variación del oponente. Medida: ESC-13. \\ \hline
\end{tabla}

\textbf{Lo que falta después de estas ocho.} Un driver con varias medidas necesita más de una iteración, y aquí solo está la primera de cada uno. DRV-01 y DRV-02 comparten la que falta más adelante: la que hace que una desconexión no ordenada no deje la partida indefinida (ESC-08), sobre la sesión de partida, que no se puede decidir antes de que exista el protocolo de la iteración~2. DRV-06 y DRV-08 tienen también trabajo pendiente: el acceso con segundo factor que un niño completa solo (ESC-03) y los otros dos puntos de variación, el extremo de red (ESC-14) y la marca con los recursos de terceros (ESC-15).

% ------------------------------------------------------------
\subsection{Por qué DRV-01 va primero}

DRV-01 obtiene la calificación máxima en los tres criterios, y es el único que la obtiene.

\begin{itemize}
  \item \textbf{Alcance estructural.} Decide quién produce el estado de la partida y quién lo consume. Cambiar de servidor autoritativo a cliente autoritativo obliga a rehacer a la vez el cliente, el servidor, el protocolo y la capa de datos: cada mensaje cambia de significado, porque deja de ser una propuesta y pasa a ser un hecho.
  \item \textbf{Costo de cambio tardío.} Es la decisión que CRN-10 describe como la más cara de cambiar, y tres ASR dependen de ella: ASR-02, ASR-04 y ASR-05 fallan a la vez si el estado vigente no vive en la memoria del servidor.
  \item \textbf{Drivers que habilita.} DRV-02 no puede fijar el apretón de manos ni la reasociación de sesión sin saber que el estado sobrevive a la conexión; DRV-03 no puede repartir el segundo sin saber cuántos viajes hay por jugada; DRV-04 no puede definir qué se escribe y cuándo sin saber quién tiene el estado vigente; y DRV-06 aplica el filtro de chat donde vive la autoridad.
\end{itemize}
